\clearpage\thispagestyle{empty}
\pdfbookmark[0]{Quellen zum Auszug über die barometrische Höhenmessung von 1806}{lindenau1806laplace-sources}
\begin{center}\textbf{Quellen zum Auszug über die barometrische Höhenmessung von 1806}\end{center}
\noindent Diese KI-Übertragung ins heutige Deutsch folgt dem Auszug auf S.~485--486 der \href{https://gdz.sub.uni-goettingen.de/id/PPN236007467}{\textit{Werke}, Band~VI} (Göttingen, 1874). Die Überschrift der Werkausgabe nennt Bernhard von Lindenau als Verfasser. Er zitiert Carl Friedrich Gauß zu der Frage, ob sich das gesuchte Element mit Rohdes Methode bestimmen lässt, und nimmt dann selbst dazu Stellung.
\par\smallskip\noindent Der Auszug stammt aus dem Brief vom 18. Mai 1806 in der \textit{Monatlichen Correspondenz}, Band~XIV, Juli 1806, \href{https://api.digitale-sammlungen.de/iiif/image/v2/bsb10538606_00066/full/1500,/0/default.jpg}{S.~56} und \href{https://api.digitale-sammlungen.de/iiif/image/v2/bsb10538606_00067/full/1500,/0/default.jpg}{S.~57}. Die Werkausgabe setzt nach den barometrischen Rechnungen ein und lässt eine Fußnote der Zeitschrift aus. Im Zitat steht in der Werkausgabe ein Punkt, wo die Zeitschrift ein Komma setzt. KI-Übertragung und KI-Satz.
\clearpage
